\providecommand{\Needspace}[1]{}
\Needspace{8\baselineskip}
\item \textbf{Señal de temperatura de un satélite en órbita}.
\nopagebreak[4]

\begin{tcolorbox}[colback=blue!2!white,colframe=blue!55!black,title={Enunciado},breakable]
Un satélite registra la temperatura de uno de sus paneles mientras orbita la Tierra. Debido a la alternancia entre zonas iluminadas y zonas de sombra, la temperatura puede modelarse como una señal periódica con varias oscilaciones sinusoidales.

Considere que $t$ representa el tiempo medido en minutos. La temperatura del panel se modelará mediante la función
\[
T(t)
=
22
+
b_1 \sin\left(\frac{2\pi t}{q}\right)
+
b_2 \cos\left(\frac{4\pi t}{q}+\alpha_1\right)
+
b_3 \sin\left(\frac{6\pi t}{q}+\alpha_2\right),
\]
donde los parámetros $b_1$, $b_2$, $b_3$, $q$, $\alpha_1$ y $\alpha_2$ son constantes.

Use los valores
\[
b_1 = 6.5,
\qquad
b_2 = 2.8,
\qquad
b_3 = 1.4,
\qquad
q = 90,
\qquad
\alpha_1 = \frac{\pi}{4},
\qquad
\alpha_2 = -\frac{\pi}{6}.
\]

\begin{itemize}
    \item[(a)] Cree un arreglo \texttt{t} con $400$ valores entre $0$ y $180$ minutos usando \texttt{np.linspace}.

    \item[(b)] Cree una función llamada \texttt{temperatura\_panel(t)} que reciba el arreglo de tiempos y retorne la temperatura $T(t)$.

    \item[(c)] Evalue la función anterior para obtener un arreglo \texttt{T} con los valores de temperatura.

    \item[(d)] Grafique la temperatura $T$ en función del tiempo $t$ usando \texttt{matplotlib}. La gráfica debe tener título, etiquetas en los ejes, grilla y una apariencia clara.

    \item[(e)] Calcule e imprimir la temperatura máxima $T_{\max}$, la temperatura mínima $T_{\min}$, la temperatura promedio $\overline{T}$ y la amplitud térmica
    \[
    A_T = \frac{T_{\max}-T_{\min}}{2}.
    \]

    \item[(f)] Cuente cuántos puntos del arreglo \texttt{T} quedan bajo la temperatura promedio $\overline{T}$.

    \item[(g)] Use un condicional para clasificar la variación térmica como moderada o intensa. Para ello, defina la variación relativa como
    \[
    R_T=\frac{A_T}{\overline{T}}.
    \]
    Si $R_T \geq 0.20$, la variación térmica se considerará intensa. En caso contrario, se considerará moderada.
\end{itemize}
\end{tcolorbox}

\begin{solution}
Una posible solución consiste en construir el arreglo temporal, definir una función vectorizada para el modelo de temperatura y luego calcular las cantidades pedidas directamente con NumPy.

\begin{pythonjupyter}
import numpy as np
import matplotlib.pyplot as plt

# Parámetros del modelo de temperatura
b1 = 6.5
b2 = 2.8
b3 = 1.4
q = 90
alpha1 = np.pi / 4
alpha2 = -np.pi / 6

# (a) Arreglo de tiempo entre 0 y 180 minutos
t = np.linspace(0, 180, 400)

# (b) Función para calcular la temperatura del panel
def temperatura_panel(t):
    T = (
        22
        + b1 * np.sin(2 * np.pi * t / q)
        + b2 * np.cos(4 * np.pi * t / q + alpha1)
        + b3 * np.sin(6 * np.pi * t / q + alpha2)
    )
    return T

# (c) Evaluar la temperatura en cada instante
T = temperatura_panel(t)

# (d) Graficar temperatura versus tiempo
plt.figure(figsize=(9, 5))
plt.plot(t, T, linewidth=2, label="Temperatura")
plt.axhline(np.mean(T), linestyle="--", linewidth=1.5, label="Promedio")

plt.title("Temperatura simulada de un panel satelital")
plt.xlabel("Tiempo t [min]")
plt.ylabel("Temperatura T(t) [C]")
plt.grid(True, alpha=0.3)
plt.legend()
plt.show()

# (e) Cálculo de cantidades características
T_max = np.max(T)
T_min = np.min(T)
T_promedio = np.mean(T)
amplitud_termica = (T_max - T_min) / 2

print(f"Temperatura máxima: {T_max:.3f}")
print(f"Temperatura mínima: {T_min:.3f}")
print(f"Temperatura promedio: {T_promedio:.3f}")
print(f"Amplitud térmica: {amplitud_termica:.3f}")

# (f) Contar puntos bajo la temperatura promedio
puntos_bajo_promedio = np.sum(T < T_promedio)

print(f"Puntos bajo el promedio: {puntos_bajo_promedio}")

# (g) Clasificación de la variación térmica
variacion_relativa = amplitud_termica / T_promedio

if variacion_relativa >= 0.20:
    clasificacion = "intensa"
else:
    clasificacion = "moderada"

print(f"Variación relativa: {variacion_relativa:.3f}")
print(f"La variación térmica es {clasificacion}.")
\end{pythonjupyter}
\end{solution}
